\section{Especificaciones}
\label{sec:especificaciones}

Para ambos tipos se exige una atenuación máxima en la banda de paso $A_p = \qty{1}{\decibel}$ y una atenuación mínima en la banda de
rechazo $A_s = \qty{24}{\decibel}$, con los bordes de la \cref{tab:especificaciones}.

\begin{table}[H]
  \centering
  \caption{Plantilla de atenuación asignada al grupo~3.}
  \label{tab:especificaciones}
  \begin{tabular}{@{}lccc@{}}
    \toprule
    Tipo & Banda(s) de paso & Banda de rechazo & $A_p$ / $A_s$ \\
    \midrule
    Pasabajos (LP)     & $f \le \qty{1000}{\hertz}$ & $f \ge \qty{4000}{\hertz}$ & \qty{1}{\decibel} / \qty{24}{\decibel} \\
    Rechazabanda (BS)  & $f \le \qty{500}{\hertz}$ \;o\; $f \ge \qty{5500}{\hertz}$ & $\qty{1200}{\hertz} \le f \le \qty{2400}{\hertz}$ & \qty{1}{\decibel} / \qty{24}{\decibel} \\
    \bottomrule
  \end{tabular}
\end{table}

\paragraph{Convenciones.} La atenuación se define como
\begin{equation}
  A(\omega) = -20\log_{10}\lvert H(\jw)\rvert \quad [\mathrm{dB}],
  \label{eq:atenuacion}
\end{equation}
referida al máximo de $\lvert H(\jw)\rvert$; es decir, las funciones se normalizan con $\max_\omega \lvert H(\jw)\rvert = 1$ (convención
de \texttt{scipy.signal}~\cite{virtanen2020}). Las frecuencias se especifican en hertz y se convierten a rad/s con $\omega = 2\pi f$ para
todo el cálculo. Una función \emph{cumple} la plantilla si $A(\omega) \le A_p$ en \emph{toda} la banda de paso y $A(\omega) \ge A_s$ en
\emph{toda} la banda de rechazo, no solo en los bordes.
